\section{Introduction}\label{sec_intro}

\subsection{Context and purpose}
%\bruno{Je ne pense pas que ce soit une super idée d'attaquer par le côté appliqué, ça lui donne trop de poids. C'est mieux d'attaquer sur le côté non-linéaire, et surtout adaptatif.} Pas Faux !

Time-frequency transforms and generalisations (wavelets and others) have long been used in various theoretical and applied domains. Besides quadratic transforms (Wigner distributions and generalizations), linear transforms such as the Gabor/STFT~\cite{daubechies1992ten,grochenig2013foundations} and wavelet transforms~\cite{Grossmann1984decomposition,daubechies1992ten,mallat2008wavelet,meyer1993wavelets} generally enjoy simple and useful invertibility properties, and therefore allow describing functions and signals as linear combination of building blocks, called time-frequency atoms. The time (and frequency/scale) resolution of the latter is given by the construction rule: constant time and frequency resolution for Gabor/STFT, generated by translations and modulations, and constant relative frequency resolution for wavelets, generated by translations and rescalings. See also~\cite{Kalisa1993weyl,Ali2000coherent,Fornasier2007banach} for alternative constructions that implement other scaling rules. Variants were also considered in specific applied domains, such as the Stockwell transform~\cite{Stockwell1996localization} in geophysics, which is very close to the constant $Q$ transform~\cite{Brown1991calculation,velasco2011constructing} we consider below, and the continuous wavelet transform.

In several application domains, in particular audio signal processing, it has been shown that adapting the scale of time-frequency atoms to the content of the signal can provide more efficient signal descriptions. The window size tunes the time-frequency resolution of the analysis. The latter is constrained by the uncertainty principle, which can be given various quantitative formulations (see~\cite{Folland1997uncertainty,Ricaud2013refined} and references therein), and which basically states that precision in time domain is possible at the price of precision loss in frequency domain, and vice versa. The problem of tuning time-frequency resolution as a function of time or frequency has been addressed by various authors, in various contexts, mostly with pre-defined dependence.  As a motivation, even though there is no complete consensus on psycho-physical aspects of human perception, it is known to involve several non-linear effects~\cite{Oxenham2018how}, and it has been claimed that this non-linearity allows going beyond time-frequency uncertainty in terms of localisation~\cite{Oppenheim2013human}. Such strategy has been successfully implemented in some advanced audio coders such as AAC (see~\cite{Brandenburg1999mp3} for a short account), which can switch dynamically between short and long local cosine windows. Another possible motivation can be  hyper-resolution (separation of close locally harmonic components in signals). However, the adaptation is often driven by prior heuristic computations, and the inversion of the corresponding transforms (if any) results from \textit{ad hoc} adaptations. To our knowledge, the non-linear problem where the time-frequency resolution is adapted to the analysed function, has not been considered from the mathematical point of view so far.

\smallskip
The goal of the present paper is to introduce and study such adaptive time-frequency transforms, able to adapt their time-frequency resolution to the analysed function.
This is done here by introducing a \textit{focus function} $f\mapsto\sigma_f$ which adapts the shape (size, bandwidth etc...) of the analyse window to specific properties of the analysed signal $f$. %, and then then allows the transform to adapt themselves to the signal. 
Stepping away from fixed time-frequency resolution complexifies significantly the analysis.
The purpose of this article is to introduce properly the non linear transforms, prove that this class of transform are well defined on $L^2$ and provide explicit, signal-dependent, lower and upper bounds for their norm (which depends on the choice of focus function).
We first introduce in Section~\ref{sec_CSQT} a generalisation of a constant Q transform studied in~\cite{velasco2011constructing,holighaus2012framework}, and prove an isometry result, using notations and ideas that will be useful for the rest of the paper. Section~\ref{sec_frequency_focus} introduces our first adaptive transform $M^\nu$ that associates with a signal $f$ a focus function $\sigma_f^\nu: \,\omega\mapsto\sigma^\nu_f(\omega)$ defined in the frequency domain. We then prove in Theorem~\ref{theorem_norm_control_frequency_focus} the well-definedness of $M^\nu$ as a map from $L^2(\R)$ into $L^2(\R^2)$ and obtain a norm control of the form $c_f\|f\|\le \|M^\nu f\|\le C_f\|f\|$, with explicit constants $c_f,C_f$, under suitable assumptions on the focus function $\sigma_f^\nu$. In section~\ref{sec_time_focus} we introduce a time-focused transform $M^\tau$ which associates with a signal $f$ a focus function of $\sigma_f^\tau:\,t\mapsto \sigma^\tau_f(t)$ defined in the time domain. We prove similar norm control and $L^2$ well-definedness results, 
%however the method is much more straightforward but we still prove it through the use of 
using a suitable non linear kernel. We also give in Sections~\ref{sec_frequency_focus} and~\ref{sec_time_focus} explicit examples of focus functions which may be of interest in applications, and illustrate their behaviour on toy examples in the Appendix. Section~\ref{sec_conclusion} is devoted to conclusions and perspectives.





\subsection{Notations}

We first introduce or recall a few notations. We will often use the notation $C([a,b],X)$ for the space of continuous functions from $\R$ into $X$, supported  in $[a,b]$. More generally, given $A\subset \R^\R$ we denote by $A_c$ the subspace of functions in $A$ with compact support, and by $A_0$ the subspace of functions in $A$ which vanish at infinity. We define $A+c:=\{f+c, f\in A\}$ for $c\in\R$ and $A^+$ denotes the subspace of non-negative valued functions on $A$.

Given an open interval $I$ in $\R$, $C^n(I)$ stands for the space of function which are $n$ times derivable on $I$, and $C^\infty(I) = \cap_{n\geqslant 0} C^n(I)$, furthermore we denote  by $C_p(I)$ the space of piecewise continuous functions on $I$. Finally, we denote by $\mathcal{E}$ the space of step functions, and by $\mathcal{E}([a,b])$ the space of step function with support included in $[a,b]$.

We recall that $L^p(\R^d)$ is the set of $p-$integrable functions (quotiented out by the classical equivalence relationship), for $p>1$. The shorter notation $L^p$ refers to $L^p(\R)$. When no measure is specified it means that we are working with the Lebesgue measure, otherwise we write $L^p(\R^d, d\mu)$ with $d\mu$ a regular enough measure on $\R^d$. For functions of several variables, we use the notation  $\R^d_t$ to specify that $t\in\R^d$. The notation $L^p(\R^d_t, L^q(\R^l_x))$ stands for functions $f$ for which the above mixed norm is finite
\begin{equation}
    \|f\|_{L^p_t L^q_x} := \left( \int_{\R^d}\left(\int_{\R^l} |f(x,t)|^q dx \right)^{p/q} dt \right)^{1/p}\ .
\end{equation}
When there is no risk of confusion, $d$ and $l$ will be omitted and we will simply write $L^p_t L^q_x$. More generally, for given two normed functional spaces $A(\R^d)$ and $B(\R^l)$ then $A_t B_x = A(\R^d_t,B(\R^l_x))$ is the set of function $f:\R^d\times\R^l\rightarrow\C$ for which the above mixed norm is finite,
\begin{equation}
    \|f\|_{A_t B_x} := \|\|f(x,t)\|_{B_x} \|_{A_t}\ .
\end{equation}

Given $f\in L^p$ its Fourier transform is either written $\F(f)$ or $\hat{f}$ and defined by the following convention if $f\in L^1$,
\begin{equation}
%    \forall \omega\in\R^d,\quad 
    \F(f)(\omega) := \int_{\R^d} f(x)e^{-2i\pi x\cdot\omega} dx\ ,
\end{equation}
for any $\omega\in\R^d$ for which the right hand side is well-defined.

